% Generated draft: requires mathematical and source review.
\documentclass[11pt]{article}
\usepackage[margin=25mm]{geometry}
\usepackage{fontspec}
\setmainfont{Latin Modern Roman}
\usepackage{amsmath,amssymb,mathtools,unicode-math}
\setmathfont{Latin Modern Math}
\usepackage{enumitem,needspace,xurl,hyperref}
\setlength{\parindent}{0pt}
\setlength{\parskip}{5pt}
\setlength{\emergencystretch}{3em}
\newcommand{\N}{\mathbb{N}}
\newcommand{\Z}{\mathbb{Z}}
\newcommand{\Q}{\mathbb{Q}}
\newcommand{\R}{\mathbb{R}}
\newcommand{\C}{\mathbb{C}}
\newcommand{\F}{\mathbb{F}}
\newcommand{\A}{\mathbb{A}}
\newcommand{\PP}{\mathbb{P}}
\newcommand{\E}{\mathbb{E}}
\newcommand{\eps}{\varepsilon}
\newcommand{\dd}{\,\mathrm{d}}
\begin{document}
\section*{UnsolvedMath ID: 20002788}
\textbf{Source problem code:} catalogue entry
\subsection{Definitions and mathematical statement}
The source asks whether beyond endoscopy can prove sharpness of the $R(\theta)$ bounds in Jiang--Liu's work. It gives neither the group, field, definition of $R(\theta)$, local normalization, nor the exact theorem whose sharpness is intended. The aid restricts to cuspidal automorphic representations $\pi$ of $\mathrm{Sp}_{2n}$ over a totally imaginary number field, with odd $n\ge5$.

In that setting, put $\theta_0=(n-1)/2$. A global Arthur parameter is written in blocks $(\tau,b)$, or $\tau\boxtimes S_b$, with $\tau$ a suitable self-dual cuspidal general-linear representation; a character block $(\chi,n)$ has a degree-one Hecke character $\chi$ and the $n$-dimensional auxiliary factor. Full parameter admissibility and the precise $R$-bound convention require the cited theorem. The aid's reported extremal dichotomy is:
\[
 \pi\text{ attains exponent }\theta_0
 \quad\Longleftrightarrow\quad
 \psi_\pi\text{ contains }(\chi,n)\text{ for some admissible }\chi.
\]
If no such character block occurs, it reports the stronger bound
\[
 R\!\left(\frac{n-1}{2}-\frac{25}{64}\right).
\]
This is a scope-specific assertion, not a statement for all reductive groups, fields, or ranks.

For a fixed admissible $\chi$, the aid describes a shifted standard-Hecke prime average $D_{\chi,X}(\pi)$ and states
\[
 \lim_{X\to\infty}D_{\chi,X}(\pi)
 =\mathbf1_{\{(\chi,n)\subset\psi_\pi\}}.
\]
The actual prime weights, shift, Hecke normalization, and treatment of bad places are not supplied, so the advertised explicit average cannot be quoted as a complete formula from this record.

The trace criterion is conditional on a justified interchange. For a test function with nonnegative spectral weights $w_f(\pi)$, for example an admissible convolution square, suppose the relevant cuspidal trace sums converge and the limit may pass through the sum. Then the reported limiting distribution is
\[
 \lim_{X\to\infty}\sum_\pi m(\pi)w_f(\pi)D_{\chi,X}(\pi)
 =\sum_{\pi:(\chi,n)\subset\psi_\pi}m(\pi)w_f(\pi)\ge0.
\]
It is nonzero exactly when an extremal-block representation detected by that test occurs. Allowing suitable detecting tests gives an existence criterion for extremizers. Positivity of the limit is not automatic positivity of every finite prime average, and pointwise convergence alone does not justify the trace interchange.

The aid therefore proposes proving sharpness by establishing a positive geometric limit for this distribution through beyond-endoscopic methods. It expressly leaves that geometric evaluation open. The parameter and spectral reductions do not themselves establish the existence of a cuspidal extremizer or sharpness of the original bound.
\subsection{Short English statement}
The supplied reduction identifies a particular extremal Arthur block and a spectral detector in a restricted symplectic setting. Sharpness would follow from a nonzero positive limiting trace after a justified interchange. The required geometric beyond-endoscopic evaluation is not provided.
\subsection{Sources}
\begin{itemize}
\item Local UnsolvedMath registry, permanent ID 20002788, source code catalogue entry; created and updated 14 August 2026. Canonical \texttt{aim-representation-theory-notes.json}, zero-based \texttt{notes[59]}, tag \texttt{problem}, item 1.
\item AIM ``Automorphic kernel functions,'' December 2015 workshop according to the supplied verification: \url{https://aimath.org/pastworkshops/automorphkernelproblems.pdf}. Wording is marked exact; original remarks and literature field are empty. It refers to Jiang--Liu without a full paper citation or definition of $R(\theta)$.
\item Research attempt 1: result \texttt{reduction}, problem status at run open, rigor label \texttt{proved}, novelty confidence low. No complete authoritative references, explicit detector formula, or geometric positive-limit theorem are supplied.
\item Ulam AI / UnsolvedMath Contributors, CC BY 4.0. Machine-generated research aid, not peer reviewed; independently verify before citation.
\end{itemize}
\end{document}
